% Quick start guide
\documentclass[envcountsect]{beamer}

\usepackage[utf8]{inputenc}
\usepackage[T2A]{fontenc}
\usepackage[english,russian]{babel}
\usepackage{graphicx}
\usepackage{amsmath}
\usepackage{float}
\usepackage{hyperref}
\usepackage{amssymb}
\usepackage{bbm}
\usepackage{amsthm}
\usepackage{indentfirst}
\usepackage{comment}
\usepackage[thinc]{esdiff}

\selectlanguage{russian}
\usetheme{Warsaw}

\newtheorem{rutheorem}{Теорема}
\newtheorem{rustatement}{Утверждение}
\newtheorem*{ruremark}{Замечание}
\newtheorem*{ruexample}{Пример}
\theoremstyle{definition}
\newtheorem{rudefinition}{Определение}[section]

% Title page details
\title{<<Условия квазимонотонности для кусочно-аффинных систем дифференциальных уравнений>>}
\subtitle{Дипломная работа}
\author{Миловидов Михаил}
\institute{МГУ имени М. В. Ломоносова}
\date{\today}

\begin{document}
	
	\begin{frame}
		% Print the title page as the first slide
		\titlepage
	\end{frame}
	
	\begin{frame}
		\begin{rudefinition}
			Система
			\begin{equation} \label{sys}
				\dot x = f(x),\, x \in \mathbb{R}^n,\, n \ge 2 
			\end{equation}
			называется квазимонотонной, если
			\[\forall i \in \overline{1,n}\ \ \forall \alpha_j \ge 0,\ j = \overline{1,n},\ j \ne i,\ \exists \tilde j:\ \alpha_{\tilde j} \ne 0 \]
			\begin{equation} \label{qm}
			\begin{split}
				f_i(x_1+\alpha_1,\dots,x_{i-1} + \alpha_{i-1},x_i,x_{i+1}+\alpha_{i+1},\dots,x_n+\alpha_n) \ge \\
				\ge f_i(x_1,\dots,x_n).
			\end{split}
			\end{equation}
			Это свойство также называется монотонностью по внедиагональным элементам, или же условием Важевского.
		\end{rudefinition}
	\end{frame}
	
	\begin{frame}
		Для квазимонотонных систем верно следующее свойство:
		\begin{rustatement}
			Пусть начальное множество \(\mathcal{X}_0\) вписано в отрезок \([\underline{x}, \overline{x}]\). Тогда множество достижимости \(x(\mathcal{X}_0,t)\) будет лежать в отрезке \([x(\underline{x}, t),\,x(\overline{x},t)]\)
		\end{rustatement}
		В случае гладких функций \(f_i(x)\) верно следующее утверждение:
		\begin{rustatement}
			Если \(f_i(x),\, i=\overline{1,n}\) --- гладкие функции и \(\diffp{f_i}{{x_j}} > 0\ \forall i \ne j\), то \(f(x)\) --- квазимонотонна.
		\end{rustatement}
	\end{frame}

	\begin{frame}{Постановка задачи}
		В данной работе мы будем исследовать условие квазимонотонности для частных видов кусочно-афинных (кусочно-линейных) функций. В результате нужно получить некоторые необходимые и достаточные условия. Рассмотрим следующие виды функции \(f(x)\):
		\begin{enumerate}
			\item \(\forall i = \overline{1,n}\ \, f_i(x) \) --- суперпозиция функций \(\min()\), \(\max()\) и арифметических операций.
			\item Функция \(f(x)\) определена на множестве \(\Omega \subset \mathbb{R}^n\), причем \(\Omega = \bigsqcup_{k = 1}^m \Omega^{(k)}\), где \(\Omega^{(k)}\) --- выпуклые многогранники, и функция \(f(x)\) на \(\Omega^{(k)}\) имеет вид \(f_i(x) = \sum_{j = 1}^{n} a_{ij}^{(k)}x_j + b_i^{(k)}\). Значения функции на границах обговариваются отдельно.
		\end{enumerate}
	\end{frame}
	
	\begin{frame}{Афинные системы}
		Рассмотрим афинную функцию вида
		\[f(x) = Ax + b,\, x \in \mathbb{R}^{n},\, A=[a_{ij}] \in \mathbb{R}^{n \times n},\, b \in \mathbb{R}^{n}.\]
		\[f_i(x) = \sum_{j=1}^{n} a_{ij}x_j + b_i.\]
		
		Для таких функций справедлив следующий критерий квазимонотонности:
		\begin{equation} \label{qmaff}
			a_{ij} \ge 0\ \forall i \neq j,\ i,j = \overline{1, n}.
		\end{equation}
	\end{frame}
	
	\begin{frame}{Кусочно-афинные системы}
		Пусть 
		\[f(x) = \begin{cases}A^{(k)}x + b^{(k)},\quad x \in \Omega^{(k)},\end{cases}\]
		где
		\[\Omega \subset \mathbb{R}^n, \Omega = \bigsqcup_{k = 1}^m \Omega^{(k)},\, x \in \mathbb{R}^{n},\, A=[a_{ij}] \in \mathbb{R}^{n \times n},\, b \in \mathbb{R}^{n}.\]
		\[f_i^{(k)}(x) = \sum_{j=1}^{n} a_{ij}^{(k)}x_j + b_i^{(k)}.\]
		Аналогично предыдущему случаю, получаем условие:
		\begin{equation} \label{qmpwaff}
			a_{ij}^{(k)} \ge 0\ \forall i \neq j,\ i,j = \overline{1, n},\ k = \overline{1, m}.
		\end{equation}
		Однако оно уже не будет достаточным для квазимонотонности функции. Нужно также проверять неубывание при прохождении границ множеств
	\end{frame}
	
	\begin{frame}{Системы, являющиеся суперпозицией функций \(\min\) и \(\max\)}
		Рассмотрим суперпозицию функций \(\min\) и \(\max\). Рассмотрим функции, у которых \(f_i\) являются суперпозициями функций \(\min\), \(\max\) а также арифметических операций \(+\) и \(-\) и операции умножения на константу.
		\begin{ruexample}
			\[\begin{cases}
				f_1(x) = 1 + \min(2x_1,x_2) + \min(x_1, 2x_2), \\
				f_2(x) = -\max(2x_1 - x_2, x_2, \min(3, x_1)).
			\end{cases}\]
		\end{ruexample}
		Такие функции будут являться непрерывными, а значит, условие \ref{qmpwaff} будет также и достаточным. 
	\end{frame}
	
	\begin{frame}{Случай простой суперпозиции}
		Рассмотрим случай простой суперпозиции:
		\[f_i(x) = \max(g_i^{(1)}(x),\dots,g_i^{(m_i)}(x))\]
		или
		\[f_i(x) = \min(g_i^{(1)}(x),\dots,g_i^{(m_i)}(x)),\]
		где
		\[g_i^{(k)}(x) = \sum_{j=1}^n a_{ij}^{(k)} x_j + b_i^{(k)},\, i = \overline{1, n},\, k = \overline{1, m_i}.\]
		Для таких функций условие \ref{qmpwaff} принимает вид
		\begin{equation} \label{qmmm}
			a_{ij}^{(k)} \ge 0\ \forall i \neq j,\ i,j = \overline{1, n},\ k = \overline{1, m_i}.
		\end{equation}
	\end{frame}
	
	\begin{frame}{Общий случай}
		К этому случаю можно свести некоторые более сложные виды функций, однако в общем случае не всегда можно утверждать об отсутсвии квазимонотонности при присутствии отрицательного внедиагонального коэффициента, так как соответствующая функция может не попасть в итоговую кусочно-афинную функцию.
		\begin{ruexample}
			\[f_1(x) = \max(0, \min(x_2, -x_2))\]
			Упростив, получаем, что \(f_1(x) \equiv 0\), однако здесь присутствует отрицательный коэффициент \(-1\).
		\end{ruexample}
	\end{frame}
	
	\begin{frame}{Функция, заданная на симплексах}
		Рассмотрим некоторое разбиение области \(\Omega\) на симплексы, пересекающиеся только по граничным точкам. Идея состоит в том, что, задав функцию лишь на всех вершинов данных симплексов, мы сможем продолжить ее на всю область.
	Для каждого симплекса \(\Omega^{(k)}\) обозначим координаты его вершин как \(g^{(k)}_1,\dots,g^{(k)}_n\). Тогда для любой точки \(x \in \Omega^{(k)}\) будет существовать единственный вектор барицентрических координат \(\alpha^{(k)} = \alpha^{(k)}(x) = (\alpha^{(k)}_1,\,\dots, \alpha^{(k)}_{n+1})\) такой, что
	\[\sum_{i = 1}^{n + 1} \alpha^{(k)}_i = 1,\quad \alpha^{(k)}_i \ge 0\ \forall i = \overline{1, n+1},\quad \sum_{i = 1}^{n+1} \alpha^{(k)}_i g^{(k)}_i = x.\]
	\end{frame}
	
	\begin{frame}
		Дополним вектор \(x\) до  \(\tilde{x} = \begin{bmatrix} x \\ 1 \end{bmatrix}\), и пусть \(\tilde{G}^{(k)} = \begin{bmatrix} g^{(k)}_{1} & \dots & g^{(k)}_{n+1} \\ 1 & \dots & 1 \end{bmatrix} \in \mathbb{R}^{(n + 1) \times (n + 1)}\). Тогда \(\tilde{G}^{(k)} \alpha^{(k)}(x) = \tilde x\). Пусть \((\tilde{G}^{(k)})^{-1} = \begin{bmatrix} H^{(k)} & h^{(k)} \end{bmatrix}\). Получим, что \(\alpha^{(k)}(x) = H^{(k)} x + h^{(k)} \).
		
		При \(x \in \Omega^{(k)}\) справедливо следующее представление функции \(f\):
		\[f(x) = F^{(k)}\alpha^{(k)}(x) = A^{(k)}x + b^{(k)},\]
		где \[F^{(k)} = \begin{bmatrix} f(g^{(k)}_{1}) & \dots & f(g^{(k)}_{n+1}) \end{bmatrix} \in \mathbb{R}^{n\times(n+1)},\]
		\[A^{(k)} = F^{(k)}H^{(k)} \in \mathbb{R}^{n\times n},\quad b^{(k)} = F^{(k)}h^{(k)} \in \mathbb{R}^{n}.\]
		Таким образом, мы получили явный кусочно-аффинный вид функции на всей области определения и можем применить \ref{qmpwaff}.
	\end{frame}
	
	\begin{frame}{Дальнейшие действия}
		\begin{itemize}
			\item Рассмотреть частные случаи симплексной сетки
			\item Рассмотреть другие способы задания кусочно линейной функции
			\item Рассмотреть случай ограниченной области (параллелепипед)
			\item Рассмотреть обобщенную квазимонотонность (неубывание и невозрастание по разным переменным)	
		\end{itemize}
	\end{frame}
\end{document}